% Next-batch embeddings: Wu InstDisc, Yu SDC, typed GPM. Not a TSS rewrite.
% Table~\ref{tab:next-batch-embeddings} is auto-generated.

\paragraph{Next-batch embeddings, existing methods.}
Batch $t$ yields a finite-dimensional object (a memory bank, a set of class
prototypes, or a representation subspace) that batch $t{+}1$ is allowed to
use. Nothing below invents a loss. TSS still sets $\eta_m$ from
$(\hat c_m,\hat\delta_m)$; these objects sit beside that map.

\emph{Nonparametric instance discrimination (Wu et al., CVPR 2018).}
On frozen FSDS blocks the encoder is $X^{(m)}$. A linear projector
$g_m:X^{(m)}\mapsto S^{k}$ is trained with a memory bank $V^{(m)}$ of one
unit vector per window. The nonparametric softmax
\begin{equation}
  P(i\mid f)
  =
  \frac{\exp(f^\top v_i/\tau)}{\sum_j\exp(f^\top v_j/\tau)}
\end{equation}
is the training criterion; $V$ is stop-grad and momentum-updated after the
step. Keys from batch $t{-}1$ are concatenated as extra negatives (He et al.,
MoCo), so the assembled embedding of this batch is the negative dictionary
of the next. Own-slot accuracy reads instance separability of $P(X^{(m)})$.
Off-diagonal cosine (collision) rises when a location hop collapses the
cloud onto one direction; clip-level text is the structural control.
This is not $P(Y\mid X)$ and is not a stepsize.

\emph{Prototypes and semantic drift (Rebuffi iCaRL; Yu et al., CVPR 2020).}
A prototype $\mu_c^{(m)}$ is the class mean of block $m$. iCaRL's NCM
classifier on batch $t{+}1$ is nearest $\mu_c$ in cosine. When $P(X)$ takes a
location hop, stale $\mu_c$ are in the wrong place. SDC estimates the
unknown drift of an old prototype from the observed drift of current-batch
points,
\begin{equation}
  \widehat{\Delta}_c
  =
  \sum_i w_{ic}\,(z_i^{(t)}-z_i^{(t-1)})
  \Big/
  \sum_i w_{ic},
  \qquad
  w_{ic}\propto\exp\bigl(-\|z_i^{(t-1)}-\mu_c\|^2/(2\sigma^2)\bigr),
\end{equation}
and writes $\hat\mu_c\leftarrow\mu_c+\widehat{\Delta}_c$ for the next NCM.
A pure mean shift is a constant vector field; SDC recovers the new class
means without old exemplars. A concept change $P(Y\mid X)$ is a different
object: interpolating current-point motion cannot undo a label permutation,
which is why $\hat\delta_m$ (not SDC) raises $\eta_m$.

\emph{GPM composed with FSDS, not rewritten.}
Saha et al.\ (ICLR 2021) store the leading SVD bases of the layer input as
the core gradient space and take the next task's step in the orthogonal
complement, $\Delta W\leftarrow(I-MM^\top)\Delta W$. On this probe the layer
input \emph{is} the FSDS block $X^{(m)}$. Residual energy
$\|X_t(I-MM^\top)\|_F^2/\|X_t\|_F^2$ is the mass of the next batch outside
the previous CGS: a covariance/support change, not Cohen's $|d|$ on the
mean. FSDS $\hat c_m$ already types the location hop; GPM types the span.
The locked signs then decide whether the projection is allowed:
\begin{equation}
  \text{project head }m
  \iff
  \hat c_m\ge\tau_c
  \text{ and }
  \hat\delta_m<\tau_d.
\end{equation}
Covariate-loud / concept-quiet: protect the old span (same sign as shrinking
$\eta_m$). Concept-loud: GPM as written would freeze a stale $P(Y\mid X)$
inside CGS, so the gate is off and TSS still raises $\eta_m$. Both quiet:
hold. Always-on GPM is the untyped ablation. Table~\ref{tab:next-batch-embeddings}
reports the three objects on the same typed streams as TSS; they are not
plugged into the TSS loop.
